\documentclass[11pt,a4paper]{article}
\usepackage[utf8]{inputenc}
\usepackage[T1]{fontenc}
\usepackage[catalan]{babel}
\usepackage{lmodern}
\usepackage[sfdefault,lining]{FiraSans}
\usepackage[a4paper,top=2.3cm,bottom=2.2cm,left=1.8cm,right=1.8cm,headheight=1.3cm,headsep=0.4cm,footskip=1cm]{geometry}
\usepackage{amsmath,amssymb}
\usepackage{icomma}
\usepackage{graphicx}
\usepackage[table]{xcolor}
\usepackage{tikz}
\usepackage{fontawesome5}
\usepackage{ccicons}
\usepackage{fancyhdr}
\usepackage{lastpage}
\usepackage{tasks}
\usepackage{needspace}
\usepackage[most]{tcolorbox}
\usepackage{array}
\usepackage[hidelinks,pdftitle={Solucions de la fitxa de la unitat 9: Derivades i aplicacions. 1r de Batxillerat MI},pdfsubject={Matemàtiques I},pdfkeywords={taxa de variació mitjana, derivada, funció derivada, regles de derivació, regla de la cadena, recta tangent, recta normal, monotonia, màxims i mínims, representació gràfica, optimització}]{hyperref}

\definecolor{navy}{HTML}{1B2748}
\definecolor{accent}{HTML}{E8603A}
\definecolor{teal}{HTML}{1F9E90}
\definecolor{soft}{HTML}{EEF1F8}
\definecolor{muted}{HTML}{5D6685}

\setlength{\parindent}{0pt}
\setlength{\parskip}{3pt}
\settasks{label=\alph*), label-format=\bfseries\color{navy}, label-width=1.3em, item-indent=1.8em, column-sep=0.6em, after-item-skip=5pt, before-skip=3pt, after-skip=4pt}

% ---------- capçalera i peu ----------
\newcommand{\llicencia}{\href{https://creativecommons.org/licenses/by-nc-sa/4.0/deed.ca}{\ccbyncsa}\ \ Llicència Creative Commons Reconeixement-NoComercial-CompartirIgual 4.0 Internacional (CC BY-NC-SA 4.0)}
\pagestyle{fancy}
\fancyhf{}
\fancyhead[L]{\raisebox{-0.35\height}{\includegraphics[height=0.95cm]{logo-gva.png}}}
\fancyhead[C]{\small\color{muted}Solucions de la fitxa · Unitat 9: Derivades i aplicacions · 1r Batxillerat MI}
\fancyhead[R]{\raisebox{-0.35\height}{\includegraphics[height=0.95cm]{logo-vermellar.png}}}
\fancyfoot[L]{\scriptsize\color{muted}\llicencia}
\fancyfoot[R]{\small\color{muted}\thepage}
\renewcommand{\headrulewidth}{0.4pt}
\renewcommand{\headrule}{\color{muted!40}\hrule height \headrulewidth}
\fancypagestyle{primera}{\fancyhf{}\renewcommand{\headrulewidth}{0pt}
  \fancyfoot[L]{\scriptsize\color{muted}\llicencia}
  \fancyfoot[R]{\small\color{muted}\thepage}}

% ---------- nivells i exercicis ----------
\newcommand{\estrelles}[1]{{\color{accent}\ifcase#1\or\faStar\faStar[regular]\faStar[regular]\or\faStar\faStar\faStar[regular]\or\faStar\faStar\faStar\fi}}
\newcounter{exer}
\newcommand{\ex}[1]{\par\needspace{4\baselineskip}\medskip\refstepcounter{exer}\noindent{\color{navy}\bfseries\theexer.}\ \ }

\newcommand{\seccio}[2]{\par\needspace{16\baselineskip}\bigskip
  \begin{tcolorbox}[colback=navy,colframe=navy,coltext=white,arc=2mm,boxrule=0pt,left=3mm,right=3mm,top=1.2mm,bottom=1.2mm]
  {\large\bfseries #1\quad #2}
  \end{tcolorbox}\vspace{-2pt}}

\newcommand{\exf}[1]{\par\medskip\refstepcounter{exer}\noindent{\color{navy}\bfseries\theexer.}\ \ }
\providecommand{\gt}{>}\providecommand{\lt}{<}
% quadre de recordatori (només a la fitxa d'exercicis)
\newcommand{\recorda}[1]{\par\smallskip\begin{tcolorbox}[colback=accent!7,colframe=accent!55,arc=2mm,boxrule=0.5pt,left=2.5mm,right=2.5mm,top=1mm,bottom=1mm]\small{\color{accent}\faLightbulb[regular]}\ \textbf{Recorda:}\ #1\end{tcolorbox}}
\newcommand{\res}[1]{\boxed{#1}}\newcommand{\rest}[1]{\fbox{#1}}

\def\fmneed{11}% que el títol d'una secció no quede sol al final del full
% =====================================================================
% Capçalera FentMates de les fitxes: un color per a cada curs.
% Els colors són els de la pàgina inicial (LEVELS a assets/js/materials.js);
% el logo ve de _eines/logo/pdf/ (python3 _eines/logo/genera.py).
%
% Ús, al final del preàmbul de la unitat:
%   % =====================================================================
% Capçalera FentMates de les fitxes: un color per a cada curs.
% Els colors són els de la pàgina inicial (LEVELS a assets/js/materials.js);
% el logo ve de _eines/logo/pdf/ (python3 _eines/logo/genera.py).
%
% Ús, al final del preàmbul de la unitat:
%   % =====================================================================
% Capçalera FentMates de les fitxes: un color per a cada curs.
% Els colors són els de la pàgina inicial (LEVELS a assets/js/materials.js);
% el logo ve de _eines/logo/pdf/ (python3 _eines/logo/genera.py).
%
% Ús, al final del preàmbul de la unitat:
%   \input{../comu/fentmates.tex}
%   \fmcurs{1eso}
% i, al començament del document:
%   \capcalera{Matemàtiques}{1r ESO}{Fitxa d'exercicis · Unitat 2}{Divisibilitat}
% Opcional abans de \input: \def\fmneed{14} (espai mínim abans de cada secció).
% =====================================================================
\makeatletter
\def\fm@logodir{../../logo/pdf}
\providecommand{\fmneed}{10}
\def\fm@cursnom{}

% clau del curs, color i color del text damunt del color
\newcommand{\fm@curs}[3]{\@namedef{fm@color@#1}{#2}\@namedef{fm@text@#1}{#3}}
\fm@curs{1eso}{E6B84A}{navy}
\fm@curs{2eso}{3B82F6}{white}
\fm@curs{3eso}{FF7A50}{white}
\fm@curs{4eso-a}{A78BFA}{white}
\fm@curs{4eso-b}{6D28D9}{white}
\fm@curs{1batx-mi}{9B1C1C}{white}
\fm@curs{1batx-ccss}{EF6B6B}{white}
\fm@curs{2batx-mii}{1E6B47}{white}
\fm@curs{2batx-ccss}{2BB3A3}{white}

\newcommand{\fmcurs}[1]{%
  \@ifundefined{fm@color@#1}{\PackageError{fentmates}{Curs desconegut: #1}{}}{}%
  \edef\fm@tmp{\noexpand\definecolor{curs}{HTML}{\@nameuse{fm@color@#1}}%
               \noexpand\colorlet{curstext}{\@nameuse{fm@text@#1}}}\fm@tmp
  \colorlet{cursfosc}{curs!72!black}%
  \def\fm@logo{\fm@logodir/vertical-#1.pdf}%
  \def\fm@marca{\fm@logodir/marca-#1.pdf}}

% ---------- capçalera de la primera pàgina ----------
% #1 matèria (petit), #2 curs, #3 tipus de document i unitat, #4 títol
\newcommand{\capcalera}[4]{%
  \gdef\fm@cursnom{#2}%
  \par\noindent
  \begin{tcolorbox}[colback=curs,colframe=curs,boxrule=0pt,arc=2.2mm,width=3.5cm,height=2.15cm,
      valign=center,left=3mm,right=2mm,top=1mm,bottom=1mm,nobeforeafter,box align=center]
    \color{curstext}{\footnotesize #1}\par\vspace{1pt}{\fontsize{19}{21}\selectfont\bfseries #2\par}%
  \end{tcolorbox}\hspace{5mm}%
  \begin{minipage}[c]{\dimexpr\linewidth-3.5cm-5mm-2.9cm\relax}
    {\color{cursfosc}\small\bfseries\MakeUppercase{#3}}\par\vspace{2pt}
    {\color{navy}\fontsize{20}{23}\selectfont\bfseries #4\par}
  \end{minipage}\hfill
  \begin{minipage}[c]{2.6cm}\raggedleft\includegraphics[width=2.6cm]{\fm@logo}\end{minipage}\par
  \vspace{2pt}}

% ---------- barres de secció del color del curs ----------
\renewcommand{\seccio}[2]{\par\needspace{\fmneed\baselineskip}\bigskip
  \begin{tcolorbox}[colback=curs,colframe=curs,coltext=curstext,arc=2mm,boxrule=0pt,left=3mm,right=3mm,top=1.2mm,bottom=1.2mm]
  {\large\bfseries #1\quad #2}
  \end{tcolorbox}\vspace{-2pt}}

% ---------- estrelles de dificultat del color del curs (a Batxillerat no se'n posen) ----------
\@ifundefined{estrelles}{}{%
  \renewcommand{\estrelles}[1]{{\color{curs}\ifcase#1\or\faStar\faStar[regular]\faStar[regular]\or\faStar\faStar\faStar[regular]\or\faStar\faStar\faStar\fi}}}

% ---------- capçalera i peu de les pàgines ----------
% primera pàgina: logos institucionals al peu
\fancypagestyle{primera}{\fancyhf{}\renewcommand{\headrulewidth}{0pt}
  \fancyfoot[L]{\raisebox{-0.32\height}{\includegraphics[height=0.78cm]{logo-gva.png}}\hspace{5mm}%
    \raisebox{-0.32\height}{\includegraphics[height=0.78cm]{logo-vermellar.png}}}
  \fancyfoot[R]{\scriptsize\color{muted}\href{https://creativecommons.org/licenses/by-nc-sa/4.0/deed.ca}{\ccbyncsa}\ \ CC BY-NC-SA 4.0 · fentmates.org\qquad\small\thepage}}
% la resta: el símbol a l'esquerra i el curs a la dreta, amb el títol al centre com abans
\fancyhead[L]{\raisebox{-0.3\height}{\includegraphics[height=0.62cm]{\fm@marca}}}
\fancyhead[R]{\small\bfseries\color{cursfosc}\fm@cursnom}
\renewcommand{\headrule}{\color{curs!45}\hrule height \headrulewidth}
\makeatother

%   \fmcurs{1eso}
% i, al començament del document:
%   \capcalera{Matemàtiques}{1r ESO}{Fitxa d'exercicis · Unitat 2}{Divisibilitat}
% Opcional abans de \input: \def\fmneed{14} (espai mínim abans de cada secció).
% =====================================================================
\makeatletter
\def\fm@logodir{../../logo/pdf}
\providecommand{\fmneed}{10}
\def\fm@cursnom{}

% clau del curs, color i color del text damunt del color
\newcommand{\fm@curs}[3]{\@namedef{fm@color@#1}{#2}\@namedef{fm@text@#1}{#3}}
\fm@curs{1eso}{E6B84A}{navy}
\fm@curs{2eso}{3B82F6}{white}
\fm@curs{3eso}{FF7A50}{white}
\fm@curs{4eso-a}{A78BFA}{white}
\fm@curs{4eso-b}{6D28D9}{white}
\fm@curs{1batx-mi}{9B1C1C}{white}
\fm@curs{1batx-ccss}{EF6B6B}{white}
\fm@curs{2batx-mii}{1E6B47}{white}
\fm@curs{2batx-ccss}{2BB3A3}{white}

\newcommand{\fmcurs}[1]{%
  \@ifundefined{fm@color@#1}{\PackageError{fentmates}{Curs desconegut: #1}{}}{}%
  \edef\fm@tmp{\noexpand\definecolor{curs}{HTML}{\@nameuse{fm@color@#1}}%
               \noexpand\colorlet{curstext}{\@nameuse{fm@text@#1}}}\fm@tmp
  \colorlet{cursfosc}{curs!72!black}%
  \def\fm@logo{\fm@logodir/vertical-#1.pdf}%
  \def\fm@marca{\fm@logodir/marca-#1.pdf}}

% ---------- capçalera de la primera pàgina ----------
% #1 matèria (petit), #2 curs, #3 tipus de document i unitat, #4 títol
\newcommand{\capcalera}[4]{%
  \gdef\fm@cursnom{#2}%
  \par\noindent
  \begin{tcolorbox}[colback=curs,colframe=curs,boxrule=0pt,arc=2.2mm,width=3.5cm,height=2.15cm,
      valign=center,left=3mm,right=2mm,top=1mm,bottom=1mm,nobeforeafter,box align=center]
    \color{curstext}{\footnotesize #1}\par\vspace{1pt}{\fontsize{19}{21}\selectfont\bfseries #2\par}%
  \end{tcolorbox}\hspace{5mm}%
  \begin{minipage}[c]{\dimexpr\linewidth-3.5cm-5mm-2.9cm\relax}
    {\color{cursfosc}\small\bfseries\MakeUppercase{#3}}\par\vspace{2pt}
    {\color{navy}\fontsize{20}{23}\selectfont\bfseries #4\par}
  \end{minipage}\hfill
  \begin{minipage}[c]{2.6cm}\raggedleft\includegraphics[width=2.6cm]{\fm@logo}\end{minipage}\par
  \vspace{2pt}}

% ---------- barres de secció del color del curs ----------
\renewcommand{\seccio}[2]{\par\needspace{\fmneed\baselineskip}\bigskip
  \begin{tcolorbox}[colback=curs,colframe=curs,coltext=curstext,arc=2mm,boxrule=0pt,left=3mm,right=3mm,top=1.2mm,bottom=1.2mm]
  {\large\bfseries #1\quad #2}
  \end{tcolorbox}\vspace{-2pt}}

% ---------- estrelles de dificultat del color del curs (a Batxillerat no se'n posen) ----------
\@ifundefined{estrelles}{}{%
  \renewcommand{\estrelles}[1]{{\color{curs}\ifcase#1\or\faStar\faStar[regular]\faStar[regular]\or\faStar\faStar\faStar[regular]\or\faStar\faStar\faStar\fi}}}

% ---------- capçalera i peu de les pàgines ----------
% primera pàgina: logos institucionals al peu
\fancypagestyle{primera}{\fancyhf{}\renewcommand{\headrulewidth}{0pt}
  \fancyfoot[L]{\raisebox{-0.32\height}{\includegraphics[height=0.78cm]{logo-gva.png}}\hspace{5mm}%
    \raisebox{-0.32\height}{\includegraphics[height=0.78cm]{logo-vermellar.png}}}
  \fancyfoot[R]{\scriptsize\color{muted}\href{https://creativecommons.org/licenses/by-nc-sa/4.0/deed.ca}{\ccbyncsa}\ \ CC BY-NC-SA 4.0 · fentmates.org\qquad\small\thepage}}
% la resta: el símbol a l'esquerra i el curs a la dreta, amb el títol al centre com abans
\fancyhead[L]{\raisebox{-0.3\height}{\includegraphics[height=0.62cm]{\fm@marca}}}
\fancyhead[R]{\small\bfseries\color{cursfosc}\fm@cursnom}
\renewcommand{\headrule}{\color{curs!45}\hrule height \headrulewidth}
\makeatother

%   \fmcurs{1eso}
% i, al començament del document:
%   \capcalera{Matemàtiques}{1r ESO}{Fitxa d'exercicis · Unitat 2}{Divisibilitat}
% Opcional abans de \input: \def\fmneed{14} (espai mínim abans de cada secció).
% =====================================================================
\makeatletter
\def\fm@logodir{../../logo/pdf}
\providecommand{\fmneed}{10}
\def\fm@cursnom{}

% clau del curs, color i color del text damunt del color
\newcommand{\fm@curs}[3]{\@namedef{fm@color@#1}{#2}\@namedef{fm@text@#1}{#3}}
\fm@curs{1eso}{E6B84A}{navy}
\fm@curs{2eso}{3B82F6}{white}
\fm@curs{3eso}{FF7A50}{white}
\fm@curs{4eso-a}{A78BFA}{white}
\fm@curs{4eso-b}{6D28D9}{white}
\fm@curs{1batx-mi}{9B1C1C}{white}
\fm@curs{1batx-ccss}{EF6B6B}{white}
\fm@curs{2batx-mii}{1E6B47}{white}
\fm@curs{2batx-ccss}{2BB3A3}{white}

\newcommand{\fmcurs}[1]{%
  \@ifundefined{fm@color@#1}{\PackageError{fentmates}{Curs desconegut: #1}{}}{}%
  \edef\fm@tmp{\noexpand\definecolor{curs}{HTML}{\@nameuse{fm@color@#1}}%
               \noexpand\colorlet{curstext}{\@nameuse{fm@text@#1}}}\fm@tmp
  \colorlet{cursfosc}{curs!72!black}%
  \def\fm@logo{\fm@logodir/vertical-#1.pdf}%
  \def\fm@marca{\fm@logodir/marca-#1.pdf}}

% ---------- capçalera de la primera pàgina ----------
% #1 matèria (petit), #2 curs, #3 tipus de document i unitat, #4 títol
\newcommand{\capcalera}[4]{%
  \gdef\fm@cursnom{#2}%
  \par\noindent
  \begin{tcolorbox}[colback=curs,colframe=curs,boxrule=0pt,arc=2.2mm,width=3.5cm,height=2.15cm,
      valign=center,left=3mm,right=2mm,top=1mm,bottom=1mm,nobeforeafter,box align=center]
    \color{curstext}{\footnotesize #1}\par\vspace{1pt}{\fontsize{19}{21}\selectfont\bfseries #2\par}%
  \end{tcolorbox}\hspace{5mm}%
  \begin{minipage}[c]{\dimexpr\linewidth-3.5cm-5mm-2.9cm\relax}
    {\color{cursfosc}\small\bfseries\MakeUppercase{#3}}\par\vspace{2pt}
    {\color{navy}\fontsize{20}{23}\selectfont\bfseries #4\par}
  \end{minipage}\hfill
  \begin{minipage}[c]{2.6cm}\raggedleft\includegraphics[width=2.6cm]{\fm@logo}\end{minipage}\par
  \vspace{2pt}}

% ---------- barres de secció del color del curs ----------
\renewcommand{\seccio}[2]{\par\needspace{\fmneed\baselineskip}\bigskip
  \begin{tcolorbox}[colback=curs,colframe=curs,coltext=curstext,arc=2mm,boxrule=0pt,left=3mm,right=3mm,top=1.2mm,bottom=1.2mm]
  {\large\bfseries #1\quad #2}
  \end{tcolorbox}\vspace{-2pt}}

% ---------- estrelles de dificultat del color del curs (a Batxillerat no se'n posen) ----------
\@ifundefined{estrelles}{}{%
  \renewcommand{\estrelles}[1]{{\color{curs}\ifcase#1\or\faStar\faStar[regular]\faStar[regular]\or\faStar\faStar\faStar[regular]\or\faStar\faStar\faStar\fi}}}

% ---------- capçalera i peu de les pàgines ----------
% primera pàgina: logos institucionals al peu
\fancypagestyle{primera}{\fancyhf{}\renewcommand{\headrulewidth}{0pt}
  \fancyfoot[L]{\raisebox{-0.32\height}{\includegraphics[height=0.78cm]{logo-gva.png}}\hspace{5mm}%
    \raisebox{-0.32\height}{\includegraphics[height=0.78cm]{logo-vermellar.png}}}
  \fancyfoot[R]{\scriptsize\color{muted}\href{https://creativecommons.org/licenses/by-nc-sa/4.0/deed.ca}{\ccbyncsa}\ \ CC BY-NC-SA 4.0 · fentmates.org\qquad\small\thepage}}
% la resta: el símbol a l'esquerra i el curs a la dreta, amb el títol al centre com abans
\fancyhead[L]{\raisebox{-0.3\height}{\includegraphics[height=0.62cm]{\fm@marca}}}
\fancyhead[R]{\small\bfseries\color{cursfosc}\fm@cursnom}
\renewcommand{\headrule}{\color{curs!45}\hrule height \headrulewidth}
\makeatother

\fmcurs{1batx-mi}

\begin{document}
\thispagestyle{primera}

\capcalera{Matemàtiques I}{1r Batx}{Solucions de la fitxa · Unitat 9}{\raggedright Derivades i aplicacions}
\vspace{8pt}

\begin{tcolorbox}[colback=soft,colframe=soft,arc=2mm,left=3mm,right=3mm,top=1.5mm,bottom=1.5mm]
\small Intenta resoldre cada exercici abans de mirar-ne la solució. Si no t'ix el mateix resultat, repassa els passos o comprova'l amb la calculadora de derivades i les eines dels apunts interactius de la unitat.
\end{tcolorbox}

% =====================================================
\seccio{1}{Taxa de variació mitjana}

\exf{1} $f(1)=5$, $f(2)=0$, $f(5)=-3$, $f(8)=12$. \ $\text{TVM}[1,2]=\frac{0-5}{1}=\res{-5}$; \ $\text{TVM}[1,5]=\frac{-3-5}{4}=\res{-2}$; \ $\text{TVM}[1,8]=\frac{12-5}{7}=\res{1}$. \\ $f(1+h)=h^2-6h+5$: \ $\text{TVM}[1,1+h]=\frac{h^2-6h+5-5}{h}=\res{h-6}$. \ Amb $h=1,4,7$: $-5$, $-2$ i $1$.

\ex{1} $\text{TVM}[0,1]=\frac{15-0}{1}=\res{15\text{ m/s}}$; \ $\text{TVM}[1,2]=\frac{20-15}{1}=\res{5\text{ m/s}}$; \ $\text{TVM}[2,4]=\frac{0-20}{2}=\res{-10\text{ m/s}}$: és negativa perquè la pilota \rest{baixa}.

\ex{1} \begin{tasks}(3)
  \task $\frac{2-1}{3}=\res{\frac13}$
  \task $\frac{8-1}{3}=\res{\frac73}$
  \task $\frac{\frac12-1}{1}=\res{-\frac12}$
\end{tasks}


% =====================================================
\seccio{2}{Derivada en un punt i funció derivada}

\exf{1} \begin{tasks}(1)
  \task $\dfrac{f(-2+h)-f(-2)}{h}=\dfrac{\frac{1}{-2+h}+\frac12}{h}=\dfrac{1}{2(-2+h)}\to\res{f'(-2)=-\frac14}$
  \task $\dfrac{-2(a+h)+4-(-2a+4)}{h}=\dfrac{-2h}{h}=-2$: \ $\res{f'(a)=-2}$ en tots els punts, perquè la gràfica és una recta de pendent $-2$.
  \task $\dfrac{f(a+h)-f(a)}{h}=6a-5+3h\to6a-5$: \ $f'(-2)=\res{-17}$, \ $f'(1)=\res{1}$, \ $f'(0)=\res{-5}$
\end{tasks}

\ex{2} \begin{tasks}(1)
  \task $\dfrac{(x+h)^2+3(x+h)-x^2-3x}{h}=\dfrac{2xh+h^2+3h}{h}=2x+h+3\to\res{2x+3}$
  \task $\dfrac{\frac{2}{x+h+1}-\frac{2}{x+1}}{h}=\dfrac{-2h}{h(x+h+1)(x+1)}\to\res{\dfrac{-2}{(x+1)^2}}$
  \task Pel conjugat: $\dfrac{\sqrt{x+h+1}-\sqrt{x+1}}{h}=\dfrac{1}{\sqrt{x+h+1}+\sqrt{x+1}}\to\res{\dfrac{1}{2\sqrt{x+1}}}$
\end{tasks}

\ex{1} $\dfrac{(3+h)^2-9}{h}=6+h\to f'(3)=6$. \ Tangent per $(3,9)$: $y=9+6(x-3)\Rightarrow\res{y=6x-9}$


% =====================================================
\seccio{3}{Regles de derivació}

\exf{1} \begin{tasks}(2)
  \task $f'(x)=\res{6x^2-10x+3}$
  \task $f'(x)=\res{\dfrac{1}{\sqrt{2x}}+\dfrac{2\sqrt[3]{5}}{3\sqrt[3]{x}}}$
  \task $f'(x)=\res{-\dfrac{3}{2x^2\sqrt x}}$
  \task $f'(x)=\res{\dfrac{x^4+3x^2}{(x^2+1)^2}}$
  \task $f'(x)=\res{2x\ln x+x}$
  \task $f'(x)=\res{\dfrac{\sqrt3}{2}\,e^x\sqrt x\,(2x+3)}$
\end{tasks}

\ex{1} \begin{tasks}(3)
  \task $\res{4x^3+\dfrac{6}{x^3}}$
  \task $\res{\dfrac{-2}{(x-1)^2}}$
  \task $\res{\sin x+x\cos x}$
  \task $\res{\dfrac{e^x(x-1)}{x^2}}$
  \task $\res{\dfrac{1}{x\ln2}+3^x\ln3}$
  \task $\res{\tan^2x}$
  \task $\res{e^x(x+1)^2}$
  \task $\res{\dfrac{1-\ln x}{x^2}}$
  \task $\res{\cos^2x-\sin^2x}$
\end{tasks}


% =====================================================
\seccio{4}{Regla de la cadena}

\exf{1} \begin{tasks}(2)
  \task $f'(x)=\res{\dfrac{x}{\sqrt{x^2+1}}}$
  \task $f'(x)=\res{e^{x^2}(1+2x^2)}$
  \task $f'(x)=\res{\dfrac{2-6x^2}{(x^2+1)^3}}$
  \task $f'(x)=\res{\dfrac{3x^2}{x^3-1}}$
  \task $f'(x)=\res{\dfrac{1}{(x^2+1)\sqrt{x^2+1}}}$
  \task $f'(x)=\res{8(2x+1)^3}$
\end{tasks}

\ex{2} \begin{tasks}(3)
  \task $\res{6x\cos(3x^2)}$
  \task $\res{-3\cos^2x\sin x}$
  \task $\res{\cos x\cdot e^{\sin x}}$
  \task $\res{\dfrac{\cos x}{\sin x}}$
  \task $\res{\dfrac{2}{\cos^2(2x+1)}}$
  \task $\res{\dfrac{1}{2x\sqrt{\ln x}}}$
  \task $\res{\dfrac{2}{x^2-1}}$
  \task $\res{\dfrac{2x}{1+x^4}}$
  \task $\res{2x\cdot2^{x^2+1}\ln2}$
\end{tasks}

\ex{2} $f'(x)=2xe^{-x}-x^2e^{-x}=e^{-x}(2x-x^2)$: \ $f'(1)=\res{\frac1e}$. \ \ $g'(x)=2\sin x\cos x$: \ $g'\!\left(\frac\pi4\right)=2\cdot\frac{\sqrt2}{2}\cdot\frac{\sqrt2}{2}=\res{1}$.


% =====================================================
\seccio{5}{Recta tangent i recta normal}

\exf{1} \begin{tasks}(1)
  \task $f(1)=2$, $f'(x)=2x+2$, $f'(1)=4$. \ Tangent: $\res{y=4x-2}$; \ normal: $y=2-\frac14(x-1)\Rightarrow\res{y=-\frac14x+\frac94}$
  \task $f(-1)=5$, $f'(x)=4x-3$, $f'(-1)=-7$. \ Tangent: $\res{y=-7x-2}$; \ normal: $y=5+\frac17(x+1)\Rightarrow\res{y=\frac17x+\frac{36}{7}}$
  \task $f(2)=3$, $f'(x)=\frac{-2}{(x-1)^2}$, $f'(2)=-2$. \ Tangent: $\res{y=-2x+7}$; \ normal: $y=3+\frac12(x-2)\Rightarrow\res{y=\frac12x+2}$
\end{tasks}

\ex{1} $f(1)=0$, $f'(1)=1$: \ $\res{y=x-1}$. \ \ $g(0)=1$, $g'(x)=2e^{2x}$, $g'(0)=2$: \ $\res{y=2x+1}$.

\ex{2} \begin{tasks}(1)
  \task $f'(x)=3x^2-6x=0\Rightarrow x=0,\ x=2$: \ $\res{(0,1)}$ i $\res{(2,\,-3)}$.
  \task $3x^2-6x=9\Rightarrow x^2-2x-3=0\Rightarrow x=3,\ x=-1$. \ En $(3,1)$: $\res{y=9x-26}$; \ en $(-1,\,-3)$: $\res{y=9x+6}$.
\end{tasks}

\ex{2} $f'(x)=2x+a$, \ $f'(1)=2+a=5\Rightarrow\res{a=3}$. \ $f(1)=4$: \ $y=4+5(x-1)\Rightarrow\res{y=5x-1}$.


% =====================================================
\seccio{6}{Monotonia i extrems}

\exf{1} \begin{tasks}(1)
  \task $f'(x)=-2x+5=0\Rightarrow x=\frac52$. \ Creixent en $\res{\left(-\infty,\frac52\right)}$ i decreixent en $\res{\left(\frac52,+\infty\right)}$. Màxim relatiu (i absolut) en $\res{\left(\frac52,\frac{21}{4}\right)}$.
  \task $f'(x)=6x^2-6x-12=6(x+1)(x-2)$. \ Creixent en $\res{(-\infty,\,-1)\cup(2,+\infty)}$, decreixent en $\res{(-1,2)}$. Màxim relatiu $\res{(-1,12)}$, mínim relatiu $\res{(2,\,-15)}$.
\end{tasks}

\ex{2} \label{ex:racional}\begin{tasks}(1)
  \task $D=\mathbb R-\{1\}$, \ $f'(x)=\frac{x^2-2x}{(x-1)^2}=0\Rightarrow x=0,\ x=2$. \ Creixent en $\res{(-\infty,0)\cup(2,+\infty)}$, decreixent en $\res{(0,1)\cup(1,2)}$. Màxim $\res{(0,0)}$, mínim $\res{(2,4)}$.
  \task $g'(x)=e^x(x+1)=0\Rightarrow x=-1$. \ Decreixent en $\res{(-\infty,\,-1)}$, creixent en $\res{(-1,+\infty)}$. Mínim $\res{\left(-1,\,-\frac1e\right)}$.
  \task $h'(x)=\frac{2x}{x^2+1}=0\Rightarrow x=0$. \ Decreixent en $\res{(-\infty,0)}$, creixent en $\res{(0,+\infty)}$. Mínim $\res{(0,0)}$.
\end{tasks}

\ex{2} $f'(x)=4x^3-16x=4x(x^2-4)=0\Rightarrow x=0,\ x=\pm2$. \ $f''(x)=12x^2-16$: \ $f''(0)=-16<0$: màxim $\res{(0,3)}$; \ $f''(\pm2)=32>0$: mínims $\res{(-2,\,-13)}$ i $\res{(2,\,-13)}$.

\ex{2} $f'(x)=3x^2-3=0\Rightarrow x=\pm1$ (tots dos dins). \ $f(-2)=0$, $f(-1)=4$, $f(1)=0$, $f(3)=20$. \ Màxim absolut $\res{(3,20)}$; mínim absolut $\res{0}$, en $x=-2$ i en $x=1$.

\ex{3} $f'(2)=12+4a=0\Rightarrow\res{a=-3}$. \ $f(2)=8-12+b=-3\Rightarrow\res{b=1}$. \ $f''(x)=6x-6$, $f''(2)=6>0$: \rest{mínim relatiu}.


% =====================================================
\seccio{7}{Representació gràfica de funcions}

\exf{2} $D=\mathbb R$; talls: $(0,0)$, $(\pm\sqrt3,0)$; \rest{senar} (simètrica respecte de l'origen); sense asímptotes (polinomi). \ $f'(x)=3x^2-3$: creix en $(-\infty,\,-1)\cup(1,+\infty)$ i decreix en $(-1,1)$; màxim $\res{(-1,2)}$ i mínim $\res{(1,\,-2)}$.\par\smallskip{\centering \begin{tikzpicture}[scale=0.45,baseline=(current bounding box.center)]\begin{scope}\draw[muted!30,very thin] (-3,-3) grid (3,3);\draw[-latex,thick] (-3,0)--(3.4,0);\draw[-latex,thick] (0,-3)--(0,3.4);\clip (-3,-3) rectangle (3,3);\draw[accent,line width=1.5pt,smooth,samples=120,domain=-2.2:2.2] plot (\x,{\x*\x*\x-3*\x});\fill[accent] (-1,2) circle (0.15);\fill[accent] (1,-2) circle (0.15);\fill[accent] (0,0) circle (0.15);\fill[accent] (-1.732,0) circle (0.15);\fill[accent] (1.732,0) circle (0.15);\node[above,font=\scriptsize] at (-1,2) {$(-1,2)$};\node[below,font=\scriptsize] at (1,-2) {$(1,-2)$};\end{scope}\end{tikzpicture}\par}

\ex{2} $D=\mathbb R-\{1\}$; tall $(0,0)$; sense simetria. \ AV $x=1$: $\lim_{x\to1^-}f=-\infty$, $\lim_{x\to1^+}f=+\infty$. \ $x^2=(x-1)(x+1)+1$: AO $\res{y=x+1}$. \ Màxim $(0,0)$ i mínim $(2,4)$ (exercici \ref{ex:racional}).\par\smallskip{\centering \begin{tikzpicture}[scale=0.36,baseline=(current bounding box.center)]\begin{scope}\draw[muted!30,very thin] (-5,-6) grid (7,10);\draw[-latex,thick] (-5,0)--(7.4,0);\draw[-latex,thick] (0,-6)--(0,10.4);\clip (-5,-6) rectangle (7,10);\draw[accent,line width=1.5pt,smooth,samples=120,domain=-5:0.8] plot (\x,{\x*\x/(\x-1)});\draw[accent,line width=1.5pt,smooth,samples=120,domain=1.15:7] plot (\x,{\x*\x/(\x-1)});\draw[teal,dashed,thick] (1,-6)--(1,10);\draw[teal,dashed,thick] (-5,-4)--(7,8);\fill[accent] (0,0) circle (0.15);\fill[accent] (2,4) circle (0.15);\node[below right,font=\scriptsize] at (2,4) {$(2,4)$};\end{scope}\end{tikzpicture}\par}

\ex{2} $D=\mathbb R$; tall $(0,0)$; \rest{senar}. \ Sense AV; AH $\res{y=0}$ ($\lim_{x\to\pm\infty}f=0$). \ $f'(x)=\frac{4(1-x^2)}{(x^2+1)^2}$: decreix en $(-\infty,\,-1)\cup(1,+\infty)$ i creix en $(-1,1)$; mínim $\res{(-1,\,-2)}$ i màxim $\res{(1,2)}$. Positiva si $x>0$.\par\smallskip{\centering \begin{tikzpicture}[scale=0.45,baseline=(current bounding box.center)]\begin{scope}\draw[muted!30,very thin] (-5,-3) grid (5,3);\draw[-latex,thick] (-5,0)--(5.4,0);\draw[-latex,thick] (0,-3)--(0,3.4);\clip (-5,-3) rectangle (5,3);\draw[accent,line width=1.5pt,smooth,samples=120,domain=-5:5] plot (\x,{4*\x/(\x*\x+1)});\fill[accent] (-1,-2) circle (0.15);\fill[accent] (1,2) circle (0.15);\node[above,font=\scriptsize] at (1,2) {$(1,2)$};\node[below,font=\scriptsize] at (-1,-2) {$(-1,-2)$};\end{scope}\end{tikzpicture}\par}

\ex{2} $D=\mathbb R$; talls $(0,0)$ i $(\pm\sqrt2,0)$; \rest{parella}; sense asímptotes. \ $f'(x)=4x^3-4x=4x(x-1)(x+1)$: decreix en $(-\infty,\,-1)\cup(0,1)$ i creix en $(-1,0)\cup(1,+\infty)$; mínims $\res{(\pm1,\,-1)}$ i màxim $\res{(0,0)}$.\par\smallskip{\centering \begin{tikzpicture}[scale=0.45,baseline=(current bounding box.center)]\begin{scope}\draw[muted!30,very thin] (-3,-2) grid (3,3);\draw[-latex,thick] (-3,0)--(3.4,0);\draw[-latex,thick] (0,-2)--(0,3.4);\clip (-3,-2) rectangle (3,3);\draw[accent,line width=1.5pt,smooth,samples=120,domain=-1.85:1.85] plot (\x,{\x*\x*\x*\x-2*\x*\x});\fill[accent] (-1,-1) circle (0.15);\fill[accent] (1,-1) circle (0.15);\fill[accent] (0,0) circle (0.15);\fill[accent] (-1.414,0) circle (0.15);\fill[accent] (1.414,0) circle (0.15);\end{scope}\end{tikzpicture}\par}


% =====================================================
\seccio{8}{Problemes d'optimització}

\exf{1} $y=20-x$, \ $P(x)=x(20-x)=20x-x^2$. \ $P'(x)=20-2x=0\Rightarrow x=10$; $P''=-2<0$: màxim. \ Els nombres són $\res{10\text{ i }10}$ i el producte, $\res{100}$.

\ex{2} Costats perpendiculars al riu $x$ i costat paral·lel $y=\frac{800}{x}$: \ $L(x)=2x+\frac{800}{x}$, $x>0$. \ $L'(x)=2-\frac{800}{x^2}=0\Rightarrow x=20$ ($L''>0$). \ $\res{20\text{ m}\times40\text{ m}}$, amb $\res{80\text{ m}}$ de tanca.

\ex{2} $V(x)=x(12-2x)^2$, $0<x<6$. \ $V'(x)=(12-2x)(12-6x)=0\Rightarrow x=2$ ($x=6$ no és vàlid). \ $\res{x=2\text{ cm}}$ i $V(2)=2\cdot8^2=\res{128\text{ cm}^3}$.

\ex{2} Si $x$ és el costat de motlura i $z$ l'altre: $C=19x+x+2z=20x+2z$ i $xz=1\,000$. \ $C(x)=20x+\frac{2\,000}{x}$, \ $C'(x)=20-\frac{2\,000}{x^2}=0\Rightarrow x=10$. \ $\res{10\text{ cm}\times100\text{ cm}}$, amb un cost de $\res{400\text{ €}}$.

\ex{3} $C(t)=(10\,000-10t^2)(2+0{,}05t)$, $t\in[0,20]$. \ $C'(t)=-1{,}5t^2-40t+500=0\Rightarrow t\approx9{,}27$. \ Amb dies enters: $C(9)=22\,515{,}5$, $C(10)=22\,500$; \ $C(0)=20\,000$, $C(20)=18\,000$. \ Li convé vendre el \rest{9é dia}: $\res{22\,515{,}50\text{ €}}$.

\ex{2} $M(x)=\frac{C(x)}{x}=x+100+\frac{3\,600}{x}$. \ $M'(x)=1-\frac{3\,600}{x^2}=0\Rightarrow x=60$. \ $\res{60\text{ unitats}}$, cost mitjà $\res{220\text{ €}}$.

\ex{3} $d^2(x)=x^2+(x^2-2)^2$ (minimitzar $d^2$ és el mateix que minimitzar $d$). \ $(d^2)'=2x+4x(x^2-2)=2x(2x^2-3)=0\Rightarrow x=0,\ x=\pm\sqrt{\tfrac32}$. \ $d^2(0)=4$, $d^2\!\left(\pm\sqrt{\tfrac32}\right)=\frac74$. \ Punts $\res{\left(\pm\frac{\sqrt6}{2},\frac32\right)}$, a distància $\frac{\sqrt7}{2}$.


\end{document}
